\documentclass[12pt]{article}

\usepackage[a4paper,margin=1in]{geometry}
\usepackage{amsmath,amssymb,amsthm}
\usepackage{tikz}
\usepackage{enumitem}
\usepackage{xcolor}
\usepackage{array}

\setlength{\parindent}{0pt}
\setlength{\parskip}{0.75em}

\newtheorem{question}{Question}
\newtheorem*{theory}{Theory to Use}
\newtheorem*{answer}{Answer}

\title{Lecture 12 Practice Questions and Step-by-Step Answers}
\author{MA4034 -- Time Series and Stochastic Processes}
\date{}

\begin{document}

\maketitle

\section*{How to Use This File}

Each question starts with the theory or equation needed. Then the answer is solved step by step.

\section{Exponential Interarrival Times and Poisson Counts}

\begin{question}
Cars arrive according to a Poisson process with rate \(6\) cars per hour.
\begin{enumerate}[label=(\alph*)]
    \item What is the probability that no cars arrive in the next \(20\) minutes?
    \item What is the probability that exactly \(3\) cars arrive in the next \(20\) minutes?
    \item What is the probability that the waiting time until the first car is more than \(20\) minutes?
\end{enumerate}
\end{question}

\begin{theory}
If \(N(t)\) is the number of arrivals in time \(t\), then
\[
N(t)\sim \operatorname{Poisson}(\lambda t),
\]
so
\[
P(N(t)=n)=e^{-\lambda t}\frac{(\lambda t)^n}{n!}.
\]

If \(T_1\) is the waiting time until the first arrival, then
\[
T_1\sim \operatorname{Exponential}(\lambda),
\]
and
\[
P(T_1>t)=e^{-\lambda t}.
\]
\end{theory}

\begin{answer}
The rate is
\[
\lambda=6 \text{ cars per hour}.
\]

The time interval is \(20\) minutes:
\[
t=\frac{20}{60}=\frac{1}{3}\text{ hour}.
\]

Therefore,
\[
\lambda t=6\cdot \frac{1}{3}=2.
\]

\textbf{(a) No cars arrive}
\[
P(N(t)=0)=e^{-2}\frac{2^0}{0!}=e^{-2}.
\]

\[
\boxed{P(N(t)=0)=e^{-2}\approx 0.1353.}
\]

\textbf{(b) Exactly 3 cars arrive}
\[
P(N(t)=3)=e^{-2}\frac{2^3}{3!}.
\]

\[
=e^{-2}\frac{8}{6}
=\frac{4}{3}e^{-2}.
\]

\[
\boxed{P(N(t)=3)=\frac{4}{3}e^{-2}\approx 0.1804.}
\]

\textbf{(c) Waiting time until first car is more than 20 minutes}

This means no car arrives during the first \(20\) minutes:
\[
P(T_1>t)=e^{-\lambda t}=e^{-2}.
\]

\[
\boxed{P(T_1>20\text{ minutes})=e^{-2}\approx 0.1353.}
\]

Notice that part (a) and part (c) have the same answer because
\[
\{T_1>t\}=\{N(t)=0\}.
\]
\end{answer}

\section{Arrival Time of the \(n\)th Event}

\begin{question}
Buses arrive according to a Poisson process with rate \(4\) buses per hour.
Let \(S_3\) be the time of the third bus.
Find
\[
P(S_3>1).
\]
\end{question}

\begin{theory}
\[
S_n=\text{time of the }n\text{th event}.
\]

The event \(S_n>t\) means that by time \(t\), fewer than \(n\) events have occurred:
\[
\{S_n>t\}=\{N(t)<n\}.
\]

Therefore,
\[
P(S_n>t)=P(N(t)<n)
=\sum_{j=0}^{n-1}e^{-\lambda t}\frac{(\lambda t)^j}{j!}.
\]
\end{theory}

\begin{answer}
Here,
\[
\lambda=4,\qquad t=1,\qquad n=3.
\]

We need
\[
P(S_3>1)=P(N(1)<3).
\]

That means
\[
N(1)=0,1,\text{ or }2.
\]

Since
\[
N(1)\sim \operatorname{Poisson}(4),
\]
we get
\[
P(S_3>1)=P(N(1)=0)+P(N(1)=1)+P(N(1)=2).
\]

\[
=e^{-4}\frac{4^0}{0!}
+e^{-4}\frac{4^1}{1!}
+e^{-4}\frac{4^2}{2!}.
\]

\[
=e^{-4}\left(1+4+8\right).
\]

\[
\boxed{P(S_3>1)=13e^{-4}\approx 0.2381.}
\]

Interpretation: There is about a \(23.81\%\) chance that the third bus has not arrived within one hour.
\end{answer}

\section{Thinning a Poisson Process}

\begin{question}
Customers arrive at a shop according to a Poisson process with rate \(12\) customers per hour.
Each customer buys a drink with probability \(0.25\), independently of other customers.
\begin{enumerate}[label=(\alph*)]
    \item What is the rate of drink-buying customers?
    \item What is the probability that exactly \(2\) drink-buying customers arrive in the next \(30\) minutes?
    \item What is the probability that no drink-buying customers arrive in the next \(30\) minutes?
\end{enumerate}
\end{question}

\begin{theory}
Thinning means keeping each event independently with probability \(p\).

If the original process has rate \(\lambda\), then the thinned process has rate
\[
\lambda p.
\]
\end{theory}

\begin{answer}
Original arrival rate:
\[
\lambda=12.
\]

Probability a customer buys a drink:
\[
p=0.25.
\]

\textbf{(a) Thinned rate}
\[
\lambda p=12(0.25)=3.
\]

\[
\boxed{\text{Drink-buying customers arrive at rate }3\text{ per hour}.}
\]

\textbf{(b) Exactly 2 drink-buying customers in 30 minutes}

The time interval is
\[
t=\frac{30}{60}=\frac{1}{2}\text{ hour}.
\]

For drink-buying customers,
\[
\lambda t=3\cdot \frac{1}{2}=1.5.
\]

Thus
\[
N(t)\sim \operatorname{Poisson}(1.5).
\]

\[
P(N(t)=2)=e^{-1.5}\frac{1.5^2}{2!}.
\]

\[
=e^{-1.5}\frac{2.25}{2}
=1.125e^{-1.5}.
\]

\[
\boxed{P(N(t)=2)=1.125e^{-1.5}\approx 0.2510.}
\]

\textbf{(c) No drink-buying customers}
\[
P(N(t)=0)=e^{-1.5}.
\]

\[
\boxed{P(N(t)=0)=e^{-1.5}\approx 0.2231.}
\]
\end{answer}

\section{Superposition of Poisson Processes}

\begin{question}
Northbound cars pass a point according to a Poisson process with rate \(3\) cars per minute.
Southbound cars pass according to an independent Poisson process with rate \(2\) cars per minute.
\begin{enumerate}[label=(\alph*)]
    \item What is the distribution of the total number of cars passing in one minute?
    \item What is the probability that exactly \(4\) cars pass in one minute?
    \item Given that exactly \(4\) cars pass in one minute, what is the probability that exactly \(3\) are northbound?
\end{enumerate}
\end{question}

\begin{theory}
Superposition means adding independent Poisson processes.

If
\[
N_1(t)\sim \operatorname{Poisson}(\lambda_1 t)
\]
and
\[
N_2(t)\sim \operatorname{Poisson}(\lambda_2 t),
\]
then
\[
N_1(t)+N_2(t)\sim \operatorname{Poisson}((\lambda_1+\lambda_2)t).
\]

Also, conditional on the total number of events, the number from one source is binomial:
\[
N_1\mid (N_1+N_2=n)\sim \operatorname{Binomial}\left(n,\frac{\lambda_1}{\lambda_1+\lambda_2}\right).
\]
\end{theory}

\begin{answer}
\textbf{(a) Total number of cars}

The total rate is
\[
\lambda_1+\lambda_2=3+2=5.
\]

For one minute,
\[
N\sim \operatorname{Poisson}(5).
\]

\[
\boxed{N\sim \operatorname{Poisson}(5).}
\]

\textbf{(b) Exactly 4 cars}
\[
P(N=4)=e^{-5}\frac{5^4}{4!}.
\]

\[
=e^{-5}\frac{625}{24}.
\]

\[
\boxed{P(N=4)=e^{-5}\frac{625}{24}\approx 0.1755.}
\]

\textbf{(c) Given 4 cars, exactly 3 are northbound}

Given total \(4\), the probability that a car is northbound is
\[
\frac{3}{3+2}=\frac{3}{5}.
\]

So
\[
X\mid N=4\sim \operatorname{Binomial}\left(4,\frac{3}{5}\right).
\]

\[
P(X=3)=\binom{4}{3}\left(\frac{3}{5}\right)^3\left(\frac{2}{5}\right)^1.
\]

\[
=4\cdot \frac{27}{125}\cdot \frac{2}{5}
=\frac{216}{625}.
\]

\[
\boxed{P(X=3\mid N=4)=\frac{216}{625}=0.3456.}
\]
\end{answer}

\section{Memoryless Property of the Exponential Distribution}

\begin{question}
The lifetime of a machine is exponential with mean \(5\) hours.
Suppose the machine has already worked for \(3\) hours.
What is the probability that it works for at least \(2\) more hours?
\end{question}

\begin{theory}
If
\[
T\sim \operatorname{Exponential}(\lambda),
\]
then
\[
P(T>s+t\mid T>s)=P(T>t).
\]

This is the memoryless property.

If the mean is \(5\), then
\[
\lambda=\frac{1}{5}.
\]
\end{theory}

\begin{answer}
We want
\[
P(T>3+2\mid T>3).
\]

By the memoryless property,
\[
P(T>5\mid T>3)=P(T>2).
\]

Since
\[
\lambda=\frac{1}{5},
\]
we get
\[
P(T>2)=e^{-\lambda\cdot 2}
=e^{-\frac{1}{5}\cdot 2}
=e^{-2/5}.
\]

\[
\boxed{P(T>5\mid T>3)=e^{-2/5}\approx 0.6703.}
\]

Interpretation: Once the machine has survived until now, the exponential model treats it like a fresh machine with respect to future waiting time.
\end{answer}

\section{Pure Birth Model}

\begin{question}
Customers arrive at a service counter according to a pure birth process with rate \(\lambda=4\) per hour.
Find the probability that exactly \(3\) customers arrive in \(30\) minutes.
\end{question}

\begin{theory}
In a pure birth process with constant birth rate \(\lambda\), the number of arrivals in time \(t\) has a Poisson distribution:
\[
p_n(t)=P(N(t)=n)=e^{-\lambda t}\frac{(\lambda t)^n}{n!}.
\]
\end{theory}

\begin{answer}
Here,
\[
\lambda=4,\qquad t=\frac{30}{60}=\frac{1}{2}.
\]

So
\[
\lambda t=4\cdot \frac{1}{2}=2.
\]

We need
\[
P(N(t)=3)=e^{-2}\frac{2^3}{3!}.
\]

\[
=e^{-2}\frac{8}{6}
=\frac{4}{3}e^{-2}.
\]

\[
\boxed{P(N(t)=3)=\frac{4}{3}e^{-2}\approx 0.1804.}
\]
\end{answer}

\section{Pure Death Model}

\begin{question}
A system starts with \(N=5\) identical components. Each component has an exponential lifetime with death rate \(\mu=0.2\) per hour.
Find the probability that exactly \(3\) components remain after \(4\) hours.
\end{question}

\begin{theory}
For a pure death model starting with \(N\) individuals, the number of deaths by time \(t\) follows a binomial form.

Each individual survives time \(t\) with probability
\[
e^{-\mu t}.
\]

Each individual dies by time \(t\) with probability
\[
1-e^{-\mu t}.
\]

Therefore, the probability that exactly \(n\) individuals remain is
\[
P(X(t)=n)=\binom{N}{n}\left(e^{-\mu t}\right)^n
\left(1-e^{-\mu t}\right)^{N-n}.
\]
\end{theory}

\begin{answer}
Here,
\[
N=5,\qquad n=3,\qquad \mu=0.2,\qquad t=4.
\]

The survival probability of one component is
\[
e^{-\mu t}=e^{-0.2(4)}=e^{-0.8}.
\]

The death probability of one component is
\[
1-e^{-0.8}.
\]

Thus
\[
P(X(4)=3)
=\binom{5}{3}(e^{-0.8})^3(1-e^{-0.8})^2.
\]

\[
\boxed{P(X(4)=3)=10e^{-2.4}(1-e^{-0.8})^2.}
\]

Numerically,
\[
\boxed{P(X(4)=3)\approx 0.3387.}
\]
\end{answer}

\section*{Essay-Type Revision Questions}

\begin{question}
Explain the difference between the exponential distribution and the Poisson distribution in the context of arrivals.
\end{question}

\begin{answer}
The exponential distribution describes waiting time.

For example, if customers arrive randomly, the exponential distribution can describe the time until the next customer arrives:
\[
T\sim \operatorname{Exponential}(\lambda).
\]

The Poisson distribution describes the number of arrivals in a fixed time interval:
\[
N(t)\sim \operatorname{Poisson}(\lambda t).
\]

So:
\[
\text{Exponential} \quad \longrightarrow \quad \text{time between events},
\]
\[
\text{Poisson} \quad \longrightarrow \quad \text{number of events in a time interval}.
\]
\end{answer}

\begin{question}
Explain why \(P(S_n>t)=P(N(t)<n)\).
\end{question}

\begin{answer}
\(S_n\) is the time at which the \(n\)th event occurs.

The event
\[
S_n>t
\]
means that the \(n\)th event has not occurred by time \(t\).

If the \(n\)th event has not occurred by time \(t\), then the number of events by time \(t\) must be less than \(n\):
\[
N(t)<n.
\]

Therefore,
\[
\boxed{P(S_n>t)=P(N(t)<n).}
\]
\end{answer}

\end{document}
